\documentclass[journal]{IEEEtran}
% ============================================================================
% arXiv PREPRINT version: standard IEEE journal (IEEEtran, two-column), SOLE
% author — matches the companion arXiv v2 (routing-oracle/paper_arxiv_v2):
% \markboth running head "Preprint, July 2026" + journal-style page numbers.
% Body matches paper_canonical (source of truth), incl. tab:sysconfig.
% Results are auto-generated from experiments/results/*.csv via make_paper_assets.py
% (tab_results.tex / tab_ablations.tex / tab_realcost.tex / tab_realverifier.tex
%  + figs/) — regenerate assets after re-running experiments; do not hand-edit numbers.
% ============================================================================
\usepackage{cite}
\usepackage{amsmath,amssymb,amsfonts}
\usepackage{algorithm}
\usepackage{algpseudocode}
\usepackage{graphicx}
\usepackage{textcomp}
\usepackage{booktabs}
\usepackage{array}

\newcommand{\Oexp}{O^{\mathrm{exp}}}
\newcommand{\Orepro}{O^{\mathrm{repro}}}
\newcommand{\bestk}{\text{best-of-}K}

\begin{document}

\title{Resample or Reroute? Budget-Aware Test-Time Model Selection for Large Language Models}

\author{Teng-Ruei~Chen%
\thanks{T.-R. Chen is with the Institute of Bioinformatics and Systems Biology,
National Yang Ming Chiao Tung University, Hsinchu 300, Taiwan, and also with
Krixvon, Taipei 100, Taiwan (e-mail: ymchen.bi04g@g2.nctu.edu.tw).}}

\markboth{Preprint,~July~2026}{Chen: Resample or Reroute?}

\maketitle

\begin{abstract}
Routing among large language models (LLMs) trades response quality against
serving cost, motivated by the reported gap between deployed routers and a
per-instance oracle. Recent analysis shows that test-time \emph{resampling}
can recover per-instance selection headroom that no single-commit router
captures; however, that guarantee holds only under an idealized oracle
equipped with correctness labels and an unconstrained budget, neither of which
a deployed system has. To the best of our knowledge, no previous work treats
resampling the committed model and rerouting to an alternative model as
competing uses of a single per-query cost budget. Therefore, this work
formulates \emph{budget-aware test-time model selection}: given a per-query
budget and an imperfect verifier, allocate each unit of budget between
resampling and rerouting so that expected correctness is maximized. An online
\emph{resample-or-reroute} (RoR) allocation policy driven by estimated
marginal correctness per unit cost is proposed, and its behavior is grounded in the recoverability asymmetry between
selection and sampling. Replay experiments on newly regenerated multi-draw correctness
tensors from an eleven-model open-weight pool over four benchmarks of
differing difficulty show that the proposed RoR policy attains a favorable
cost--quality Pareto front relative to single-route, one-commit-router,
budget-aware best-of-$K$, cascade, and random-allocation baselines for the
tested pools,
with the largest gains on the most heterogeneous benchmark; an ablation
further shows the gains are verifier-gated, shrinking as verifier quality
degrades, and robustness replays under a provider price vector and a
label-free agreement verifier delineate where the conclusions carry over.
\end{abstract}

\begin{IEEEkeywords}
Best-of-$K$ sampling, cost-aware inference, inference efficiency, large
language models, model routing, model selection, test-time compute.
\end{IEEEkeywords}


\section{Introduction}\label{sec:introduction}
\IEEEPARstart{S}{erving} a query with large language models (LLMs) increasingly
means choosing \emph{how} to spend inference compute, not just \emph{which}
model to call: inference-time cost and energy are first-class deployment
constraints~\cite{argerich2024energy}, cost- and QoS-aware resource management
for ML services is an established concern of serving
infrastructures~\cite{khan2022mlcentric}, and the space of candidate models
keeps broadening~\cite{shao2024surveyllm}. Model \emph{routing} promises to cut cost by sending each query
to the cheapest model that can answer it, motivated by the large reported gap
between a deployed router and a per-instance oracle that, in hindsight, always
picks a correct model~\cite{routerbench,llmrouterbench,frugalgpt,routellm}.

Two facts complicate this picture. First, under stochastic decoding the
per-instance oracle is not a reproducible property: it is built from single
draws, so part of the reported gap is single-draw label noise that
\emph{no} single-commit router can capture, while the rest is genuine,
recoverable specialist advantage~\cite{routinggapreal}. Second, that
recoverable component can also be reached without any router at all: test-time
\emph{resampling} (best-of-$K$ on one committed model) provably recovers the
selection floor at the oracle's own budget~\cite{routinggapreal}. The catch is
that this guarantee assumes access to correctness labels (a perfect verifier)
and spends the oracle's budget; a real serving system has an \emph{imperfect}
verifier and a \emph{fixed} cost budget per query.

This work addresses the operational question these results leave open:
\emph{given a per-query cost budget and an imperfect verifier, should a system
resample the model it already committed to, or reroute to a different
(possibly more expensive) model?} The question is cast as a budgeted
correctness-maximization problem, and an online policy is proposed that
allocates each unit of budget to whichever action---one more sample of a
candidate, or the first sample of a new candidate---has the highest estimated
marginal correctness per unit cost.

The main contributions of this work are summarized as follows.
\begin{itemize}
  \item To the best of our knowledge, this is the first work to treat
        resampling the committed model and rerouting to an alternative as
        competing uses of a single per-query budget. This \emph{budget-aware
        test-time model selection} problem is formalized as maximizing
        expected correctness under a per-query cost budget with an imperfect
        verifier (Section~\ref{sec:problem}).
  \item An online \emph{resample-or-reroute} (RoR) allocation policy (a
        marginal-gain greedy rule and a UCB variant) is proposed, together
        with an oracle allocation that serves as a reference ceiling
        (Section~\ref{sec:method}).
  \item The policy's behavior is connected to the recoverability asymmetry
        between selection and sampling established in the companion
        analysis~\cite{routinggapreal} (Section~\ref{sec:theory}).
  \item A system-level replay evaluation on four regenerated open-model pools
        reports the cost--quality Pareto front against single-route,
        one-commit-router, budget-aware best-of-$K$, cascade, and
        random-allocation baselines, together with a
        verifier-quality ablation (Sections~\ref{sec:setup}--\ref{sec:results}).
\end{itemize}

\textbf{Relation to the companion paper.} This work is a method-oriented
follow-up to the theoretical study of the router-to-oracle
gap~\cite{routinggapreal}; the differences are threefold: 1)~for the type of
contribution, that work proves a decomposition and a recoverability asymmetry,
whereas this work builds and evaluates a deployable policy; 2)~for the
information assumptions, that work grants the recovery procedure oracle
correctness labels at the oracle's own budget, whereas this work assumes only
an imperfect verifier under a fixed per-query budget; and 3)~for the role of
routing, that work adds no router by design, whereas this work actively
decides between resampling and rerouting.

The rest of this paper is organized as follows.
Section~\ref{sec:related} reviews related work.
Section~\ref{sec:problem} formulates the problem, and
Section~\ref{sec:method} presents the resample-or-reroute policy.
Section~\ref{sec:theory} grounds the policy theoretically.
Sections~\ref{sec:setup} and~\ref{sec:results} report the experimental setup
and results, Section~\ref{sec:discussion} discusses limitations, and
Section~\ref{sec:conclusion} concludes this work.

\section{Related Work}\label{sec:related}
\textbf{Cost-aware routing and cascades.} FrugalGPT~\cite{frugalgpt} and model
cascades escalate from cheaper to more expensive models under a confidence
gate; learned routers such as RouteLLM~\cite{routellm} predict, per query, a
single model to commit to; and Sheokand \emph{et al.}~\cite{sheokand2025brokered}
escalate unresolved reasoning traces from a base model to a brokered layer of
expert LLMs, a reflection-style variant of cascading. Budgeted cascades
predate LLMs: Xu \emph{et al.}~\cite{xu2014classifier} learn classifier
cascades that trade test-time evaluation cost against accuracy, and the
reject-option literature~\cite{bartlett2008reject,hendrickx2024reject}
formalizes paying a fixed cost to abstain rather than accept a likely-wrong
prediction --- the primitive behind escalation. The same accept-or-escalate structure
predates LLMs in two adjacent literatures: QoS-aware service selection chooses,
per request, which functionally equivalent service to invoke under quality and
cost constraints~\cite{zheng2011qosaware}, and adaptive edge offloading uses a
calibrated early-exit confidence to decide whether to accept a cheap local
prediction or escalate to a larger remote model~\cite{pacheco2023calibration},
a decision studied broadly in the computation-offloading
literature~\cite{mustafa2024dnnoffloading,gauttam2026edgeai}. These approaches
optimize \emph{which} model to call
but treat each call as a single draw. The proposed policy adds an orthogonal
axis: spending budget on \emph{repeated draws} of the committed model, and
choosing between that and rerouting.

\textbf{Test-time sampling.} Best-of-$N$, self-consistency~\cite{selfconsistency},
and test-time compute scaling~\cite{scalingtesttime} improve quality by drawing
multiple samples and selecting among them; Farr
\emph{et al.}~\cite{farr2025ensemble} aggregate ensemble disagreement into
uncertainty estimates that flag unreliable LLM outputs, a verifier-style
signal of the kind the proposed policy consumes. These works fix the model and
scale samples; this work integrates sample scaling and model choice into one budgeted
decision, and uses the recoverability asymmetry~\cite{routinggapreal} to explain
\emph{when} resampling beats rerouting.

\textbf{Routing benchmarks.} RouterBench~\cite{routerbench},
RouterEval~\cite{routereval}, and LLMRouterBench~\cite{llmrouterbench} provide
per-(query,\,model) correctness matrices we build on; we reuse and extend
multi-draw regenerations of these pools. Reliability pitfalls of single-run
evaluation are increasingly documented in the evaluation
literature~\cite{sheikhi2026beyond}, including large score variance across
prompt variants~\cite{mizrahi2024state} --- the same concern the multi-draw
protocol addresses on the decoding axis.

\textbf{Efficient LLM inference.} An orthogonal family of levers reduces the
cost of each call rather than the allocation of calls: quantization,
compression, and related techniques are surveyed in~\cite{mussa2025efficient},
and inference-time energy profiling guides deployment
choices~\cite{argerich2024energy}. A second orthogonal family
operates at the serving-system layer: SLO-aware scheduling and provisioning of
inference workloads~\cite{hu2026brownoutserve,xu2023igniter} decides
\emph{how} to run the chosen model, whereas the proposed policy decides
\emph{which} model to sample next. Both families compose with the proposed
policy, which reallocates a fixed per-query budget across models and draws.

Table~\ref{tab:positioning} positions this work against the related families:
prior approaches either choose among models or spend repeated draws, but none
treats the two as competing uses of one explicit per-query budget under an
imperfect verifier.

\begin{table*}[t]\centering\footnotesize
\caption{Positioning against related approach families.}
\label{tab:positioning}
\begin{tabular}{lcccc}\toprule
 & Chooses & Repeated & Per-query & Imperfect \\
Approach family & among models & draws & budget & verifier \\ \midrule
Learned routers~\cite{routellm} & \checkmark & -- & -- & -- \\
Cascades~\cite{frugalgpt,xu2014classifier} & \checkmark (sequential) & -- & cost gate & confidence gate \\
Best-of-$N$ / self-consistency~\cite{selfconsistency,scalingtesttime} & -- & \checkmark & fixed $N$ & selector needed \\
LLM ensembles~\cite{farr2025ensemble} & \checkmark (parallel) & -- & -- & UQ signal \\
Reject option~\cite{bartlett2008reject,hendrickx2024reject} & abstain/accept & -- & reject cost & confidence \\
\textbf{RoR (this work)} & \checkmark & \checkmark & \checkmark & \checkmark (parametric) \\
\bottomrule\end{tabular}\end{table*}

\textbf{Best-arm identification.} Allocating draws to candidates connects to
pure-exploration bandits~\cite{audibert2010,kaufmann2016}; this work borrows their
marginal-value intuition but targets a budgeted correctness objective rather than
identifying a single best arm.

\section{Problem Formulation}\label{sec:problem}
Consider a query $i$, a pool of models $\mathcal{M}=\{1,\dots,M\}$ with per-draw
costs $c_1,\dots,c_M$, and a per-query cost budget $B$. Drawing a sample from
model $m$ yields a candidate answer that is correct with (unknown) probability
$p_{im}$. A (possibly imperfect) verifier $v$ scores candidates and is used both
to \emph{select} a final answer and to \emph{guide} allocation; $v$ is not
assumed to reveal ground-truth correctness. Table~\ref{tab:notation} summarizes
the notation.

\begin{table}[t]\centering
\caption{Main notation.}
\label{tab:notation}
\resizebox{\columnwidth}{!}{%
\begin{tabular}{ll}\toprule
Symbol & Meaning \\ \midrule
$i$ & query index \\
$\mathcal{M}$, $M$ & model pool and its size \\
$c_m$ & cost of one draw from model $m$ \\
$B$ & per-query cost budget \\
$p_{im}$ & probability that one draw of $m$ answers $i$ correctly \\
$k$ & recorded draws per (query, model) cell in the replay \\
$v$, $q$ & verifier and its quality parameter \\
$\bar{p}_m$ & offline (train-split) accuracy of model $m$ \\
$n_{im}$, $w_{im}$ & draws of $m$ spent on query $i$, and how many verified correct \\
$\hat{p}_{im}$ & posterior mean of $p_{im}$ during a query \\
$s$ & prior pseudo-count in the belief update \\
$\hat{a}$ & final answer returned by the policy \\
\bottomrule\end{tabular}}\end{table}

A policy $\pi$ observes the history of drawn candidates and their verifier
scores and, at each step, chooses an action $a\in\{\textsc{resample}(m),
\textsc{reroute}(m')\}$ (one more draw from an already-used model $m$, or the
first draw from a new model $m'$) until the spent budget reaches $B$; it then
returns a final answer $\hat{a}$. The objective is
\begin{equation}
\max_{\pi}\; \mathbb{E}\big[\,\mathbf{1}\{\hat{a}\ \text{correct}\}\,\big]
\quad\text{s.t.}\quad \textstyle\sum \text{cost} \le B .
\end{equation}
We report performance as a cost--quality curve: expected correctness as a
function of the average budget $B$, and compare policies by Pareto dominance.
Note that the constraint binds per query, not in expectation over queries;
with a verifier that stops a query early, the \emph{spent} cost can therefore
be far below $B$, which is exactly the effect the cost axis of the Pareto
curves measures.

\section{Method: Resample-or-Reroute}\label{sec:method}
The proposed resample-or-reroute (RoR) policy maintains, for each candidate
model, a running estimate of its success probability on the \emph{current}
query and its per-draw cost, and repeatedly spends the next unit of budget
where the estimated return per unit cost is highest. This section presents the
belief update, the allocation rule, the two variants, and the computational
cost.

\subsection{Per-Query Belief Update}\label{sec:belief}
Let $\bar{p}_m$ denote model $m$'s offline accuracy, calibrated once on the
train split. During a query $i$, after $n_{im}$ draws of model $m$ of which
$w_{im}$ were verified correct, the policy scores $m$ by the posterior mean
\begin{equation}\label{eq:belief}
\hat{p}_{im}=\frac{s\,\bar{p}_m+w_{im}}{s+n_{im}},
\end{equation}
where the pseudo-count $s$ controls how strongly the offline prior anchors the
estimate. The choice of $s$ implements the resample-or-reroute trade-off: with
a small $s$, one or two failed draws pull $\hat{p}_{im}$ below a fresh model's
$\bar{p}_{m'}$ and thereby trigger a reroute, whereas a run of successes keeps
the policy resampling the committed model. For a model not yet tried on this
query, $n_{im}=0$ and $\hat{p}_{im}=\bar{p}_m$, so rerouting competes on the
offline prior alone. Section~\ref{sec:results} shows the policy is insensitive
to $s$ over an order of magnitude.

\subsection{Marginal-Gain Allocation Rule}
At each step the policy takes the action with the highest estimated
\emph{marginal correctness per unit cost}:
\begin{equation}\label{eq:rule}
a^\star=\arg\max_{a}\ \frac{\widehat{\Delta\mathrm{corr}}(a)}{\mathrm{cost}(a)},
\end{equation}
where $\widehat{\Delta\mathrm{corr}}$ estimates the increase in the probability
that the finally selected answer is correct; under early stopping this is the
probability $\hat{p}_{im}$ that the next draw of $m$ succeeds, so the rule
reduces to $a^\star=\arg\max_{m\ \text{affordable}}\ \hat{p}_{im}/c_m$.
Resampling model $m$ increases the chance that a verifier-selected answer from
$m$ is correct (a $\bestk$-style lift); rerouting to $m'$ opens a new, possibly
higher-$p$ candidate at a different cost. Both actions are scored on the same
scale, which is what makes them competing uses of one budget rather than
separate mechanisms.

\subsection{Variants and Reference Ceiling}
We consider a greedy variant and a UCB variant that adds an exploration bonus
$\sqrt{2\ln(t+1)/(n_{im}+1)}/c_m$ at step $t$ to under-sampled models; an
oracle allocation (with access to true $p_{im}$) provides a non-deployable
upper bound. Algorithm~\ref{alg:ror} summarizes the greedy variant.

\subsection{Computational Cost}
Each step scores at most $M$ actions with $O(1)$ arithmetic per action, and a
query takes at most $\lceil B/c_{\min}\rceil$ steps, so the per-query overhead
is $O(M\,B/c_{\min})$ time and $O(M)$ memory --- negligible relative to a
single LLM call. The policy is stateless across queries except for the $M$
offline priors $\bar{p}_m$.

\begin{algorithm}[t]
\caption{Resample-or-Reroute (greedy)}
\label{alg:ror}
\begin{algorithmic}[1]
\State \textbf{Input:} pool $\mathcal{M}$, costs $c_m$, budget $B$, verifier $v$
\State spent $\gets 0$; draw one sample from the cheapest model; init $\hat{p}$
\While{spent $< B$}
  \For{each action $a$ (resample used $m$, or reroute to new $m'$)}
    \State estimate marginal gain $\widehat{\Delta\mathrm{corr}}(a)$ and cost
  \EndFor
  \State take $a^\star=\arg\max_a \widehat{\Delta\mathrm{corr}}(a)/\mathrm{cost}(a)$
  \State update $\hat{p}$; spent $\gets$ spent $+\ \mathrm{cost}(a^\star)$
\EndWhile
\State \Return answer maximizing verifier score across all drawn candidates
\end{algorithmic}
\end{algorithm}

\section{Theoretical Grounding}\label{sec:theory}
% Keep light; the heavy theory lives in the companion paper [routinggapreal].
The recoverability asymmetry~\cite{routinggapreal} states that the per-instance
selection floor is closed by no single-commit router but is recovered by
resampling at the oracle's budget. The proposed policy inherits two consequences.
First, when a committed model's reproducible success probability is high,
resampling has larger marginal correctness-per-cost than rerouting, so the
greedy rule prefers resampling---matching the asymmetry. Second, with an
imperfect verifier the reachable correctness interpolates between the
verifier-free aggregate and the labeled oracle bound
($\Orepro\le O^{\mathrm{agg}}\le\Oexp$ in the companion paper's notation); we
characterize this dependence empirically in Section~\ref{sec:results}.

\section{Experimental Setup}\label{sec:setup}
\textbf{Data.} We evaluate on newly regenerated multi-draw correctness tensors built
with the companion protocol~\cite{routinggapreal}: for each of four benchmarks
--- GSM8K (arithmetic, near-saturated), MATH-500 (competition math,
intermediate), GPQA-Diamond (graduate science, hard and heterogeneous), and
HumanEval+ (code generation, execution-scored) ---
every (query, model) cell holds $k{=}30$ seed-aligned draws at $T{=}0.2$
(top-$p$ $1.0$) from a pool of $11$ open-weight models spanning eight
pretraining lineages (Mistral; Qwen~2.5 at 7B/14B/32B plus a Qwen-based
DeepSeek-R1 distill; Phi-4; OLMo-2; Yi-1.5; Granite-3.3; Gemma-2; Llama-3.1).
The first three benchmarks are exact-match scored; on HumanEval+ a draw is
correct iff it passes the full EvalPlus test suite (per-draw pass@1). Sizes:
$500$/$500$/$198$/$164$ queries.

\textbf{Protocol.} Policies are replayed offline on the recorded draws: a
``sample from model $m$'' consumes one of its $k$ pre-computed draws (without
replacement, order randomized per trial). Queries are split 50/50; per-model
priors and the best single model are calibrated on the train half; we report
test-half accuracy averaged over $20$ random draw orderings, sweeping the
per-query budget $B$. The replay is implemented in Python/NumPy and runs on
CPU only, so every reported number is reproducible from the released tensors
without model inference.

Table~\ref{tab:settings} lists the experimental settings; parameters follow the
companion protocol~\cite{routinggapreal}, and the budget grid spans from one
draw of the cheapest model to six draws of the most expensive one.

\begin{table}[t]\centering
\caption{Experimental settings.}
\label{tab:settings}
\resizebox{\columnwidth}{!}{%
\begin{tabular}{ll}\toprule
Setting & Value \\ \midrule
Benchmarks (queries) & GSM8K (500), MATH-500 (500), GPQA-Diamond (198), HumanEval+ (164) \\
Model pool & $M=11$ open-weight models, eight lineages \\
Draws per cell & $k=30$, seed-aligned \\
Decoding & $T=0.2$, top-$p$ $1.0$ \\
Train/test split & $50/50$ (priors on train, results on test) \\
Draw orderings & $20$ independent permutations per policy \\
Budget grid & $12$ points in $[c_{\min},\, 6\,c_{\max}]$ \\
Prior pseudo-count & $s=2$ (sensitivity in Table~\ref{tab:ablations}) \\
Implementation & Python/NumPy replay, CPU only \\
\bottomrule\end{tabular}}\end{table}

\textbf{Generation environment.} The correctness tensors were produced by
serving the eleven-model pool under vLLM at $T{=}0.2$ on the hardware and
NVIDIA CUDA software in Table~\ref{tab:sysconfig} --- the vLLM/PyTorch build
rests on NVIDIA's CUDA math and communication libraries (cuBLAS, cuDNN, NCCL)
alongside CUTLASS and FlashInfer attention kernels; following the companion
protocol's small-disk design, one model is resident at a time and its weights
are evicted before the next is loaded. This concerns only how the tensors were
\emph{generated}: the allocation replay analyzed in this paper is CPU-only
(Table~\ref{tab:settings}) and reproduces every reported number from the
released tensors without any model inference or accelerator.

\begin{table}[t]\centering
\caption{Generation environment (tensor production only; the allocation replay
of this paper is CPU-only). Values are the audited runtime recorded by an
automatic environment probe.}
\label{tab:sysconfig}
\resizebox{\columnwidth}{!}{%
\begin{tabular}{ll}\toprule
Component & Value \\ \midrule
CPU & AMD EPYC~7J13, $24$ vCPUs \\
RAM & $64$\,GB \\
GPU & $2\times$ NVIDIA GeForce RTX~4090, $24$\,GB each (Ada Lovelace, cc~$8.9$) \\
OS & Ubuntu~$24.04$ LTS (kernel~$6.8$) \\
NVIDIA driver & $580.159.03$ \\
CUDA toolkit / runtime & $13.0$ (NVCC~$13.2$, NVRTC~$13.0$) \\
NVIDIA math libraries & cuBLAS~$13.1$, cuDNN~$9.19$, cuSPARSELt~$0.8$ \\
NVIDIA communication & NCCL~$2.28.9$, NVSHMEM~$3.4.5$ \\
CUDA attention kernels & CUTLASS~$4.5$, FlashInfer~$0.6.12$ (vLLM backend) \\
Framework & PyTorch~$2.11.0$, vLLM~$0.23.0$, Transformers~$5.12.1$ \\
Serving & one model resident at a time (weights evicted between models) \\
\bottomrule\end{tabular}}%
\\[2pt]{\footnotesize Also installed as CUDA~13 build dependencies (pulled via
\texttt{pip} by the PyTorch/vLLM build, beyond those above): cuFFT, cuRAND,
cuSOLVER, cuSPARSE, cuPTI, nvJitLink, NVVM, NVTX, cuFile, and nvidia-ml-py; the
complete environment probe ships with the released code.}\end{table}

\textbf{Cost.} Per-draw cost is proxied by the model's parameter count in
billions --- a monotone stand-in for \$-per-token serving cost. Replacing it
with a measured price/latency vector is a deployment-specific calibration that
does not change the mechanism (Section~\ref{sec:discussion}).

\textbf{Verifier.} The main experiment assumes a reliable verifier: a policy
stops as soon as a drawn answer is verified correct (early stopping), which is
the setting where the recoverability asymmetry is fully available. Because raw
language-model confidence is known to be poorly calibrated on question
answering~\cite{jiang2021calibration}, a deployed verifier is imperfect; we
therefore degrade the verifier parametrically: with quality $q$, final selection succeeds
with probability $q\cdot\mathbf{1}\{\text{any drawn sample correct}\} +
(1-q)\cdot(\text{fraction of drawn samples correct})$, interpolating between a
perfect verifier ($q{=}1$) and a random pick over drawn samples ($q{=}0$), with
no early stopping for $q<1$.

\textbf{Baselines.} (i) \emph{single-route}: one draw of the most accurate
model; (ii) \emph{router (one commit)}: a learned per-query router's committed
model, one draw; (iii) \emph{budget-aware best-of-$K$}: the strongest
single-model baseline --- it picks, per budget, the model maximizing expected
best-of-$K$ accuracy and only resamples it; (iv) \emph{FrugalGPT-style
cascade}: escalate cheap$\to$expensive, one draw each; (v) \emph{random
allocation}; and (vi) \emph{oracle allocation}, a non-deployable ceiling that
routes each query to the cheapest model with a correct draw.

\section{Results}\label{sec:results}
% AUTO-GENERATED by experiments/make_paper_assets.py -- do not hand-edit.
\begin{table*}[t]\centering\footnotesize
\caption{Accuracy at a matched mid budget (point nearest mean cost $26$; cost = parameter-size proxy). Single-commit baselines are budget-independent and shown at their fixed operating point. Oracle allocation is a ceiling, not a deployable policy.}
\label{tab:results}
\begin{tabular}{lcccccccc}\toprule
 & \multicolumn{2}{c}{GSM8K} & \multicolumn{2}{c}{MATH-500} & \multicolumn{2}{c}{GPQA} & \multicolumn{2}{c}{HumanEval+} \\
Policy & cost & acc & cost & acc & cost & acc & cost & acc \\ \midrule
\textbf{Resample-or-Reroute (ours)} & \textbf{9.2} & \textbf{0.993} & \textbf{26.4} & \textbf{0.887} & \textbf{25.9} & \textbf{0.968} & \textbf{20.5} & \textbf{0.952} \\
ours (UCB variant) & 9.3 & 0.993 & 26.2 & 0.876 & 26.2 & 0.934 & 22.6 & 0.954 \\
budget-aware best-of-$K$ & 13.9 & 0.983 & 25.5 & 0.867 & 27.0 & 0.861 & 26.5 & 0.852 \\
cascade (FrugalGPT-style) & 12.1 & 0.992 & 25.8 & 0.847 & 26.2 & 0.926 & 22.8 & 0.952 \\
random allocation & 14.3 & 0.992 & 26.6 & 0.846 & 24.3 & 0.706 & 25.8 & 0.959 \\ \midrule
router (one commit) & 32.0 & 0.976 & 32.0 & 0.784 & 8.0 & 0.566 & 32.0 & 0.817 \\
single-route (best model) & 32.0 & 0.966 & 32.0 & 0.776 & 8.0 & 0.551 & 32.0 & 0.858 \\
oracle allocation (ceiling) & 7.0 & 1.000 & 6.7 & 0.944 & 7.0 & 1.000 & 7.1 & 0.988 \\
\bottomrule\end{tabular}\end{table*}


\begin{figure*}[t]\centering
\includegraphics[width=\textwidth]{figs/fig_pareto.pdf}
\caption{Cost--quality Pareto fronts on the four regenerated pools (verifier
$q{=}1$; mean over 20 draw orderings). The policy (blue) dominates or matches
every budget-scalable baseline across budgets on GSM8K, MATH-500, and GPQA; on
the near-saturated HumanEval+ it leads at low budget but converges with cascade
and random allocation at high budget (see text); the
amber dashed line is the non-deployable oracle-allocation ceiling.
Single-commit policies (star/plus) are budget-independent points ---
Pareto-incomparable at the lowest GPQA budgets, overtaken beyond
cost~${\sim}11$.}
\label{fig:pareto}
\end{figure*}

\textbf{Regimes across benchmarks.} From Table~\ref{tab:results} and
Fig.~\ref{fig:pareto}, the same mechanism pays off differently by regime.
On \emph{saturated} GSM8K every budget-using policy approaches the ceiling
(cascade and random reach ${\ge}0.992$; best-of-$K$ plateaus at $0.983$), so
the margin is mostly cost: the proposed policy reaches $0.993$ at mean cost
$9.2$, i.e.\ $24$--$34\%$ cheaper than the cascade ($12.1$) and budget-aware
best-of-$K$ ($13.9$), and $3.5\times$ cheaper than single-routing the best
model ($9.2$ vs $32.0$) while also being $2.7$~points more accurate ($0.993$ vs $0.966$). This
is reasonable because on a near-saturated benchmark almost every model can
eventually produce a correct draw, so the advantage comes from stopping early
on cheap models rather than from finding a specialist. On \emph{intermediate}
MATH-500 the policy is $2.1$~points above the strongest baseline at matched
cost ($0.887$ vs $0.867$) and $10$~points above the one-commit router at
$18\%$ lower cost. On \emph{hard, heterogeneous} GPQA-Diamond --- where the
pool's specialists genuinely differ --- rerouting matters most: $0.968$ vs
$0.861$ for best-of-$K$ at matched cost ($+10.7$~points), $+4.1$ over the
cascade, and roughly $+40$~points over either single-commit policy. Note one
exception: at the very lowest GPQA budgets (${\sim}8$, below one draw of most
strong models) the single-commit points ($0.55$--$0.57$ at cost $8$) beat the
policy's cheapest-first exploration ($0.35$ at cost $7$); the policy overtakes
them once the budget affords a second draw (${\sim}11$). The overall ordering
matches the companion analysis: the more heterogeneous the pool, the more of
the gap is recoverable by moving budget \emph{between} models rather than only
resampling one.

On \emph{code} (HumanEval+) the picture resembles saturated GSM8K: at matched
mid budget RoR reaches $0.952$, against $0.852$ for budget-aware best-of-$K$
($+10.0$~points) and $0.858$ for single-routing the best model at lower cost.
Here, though, the honesty cuts the other way. With a near-perfect execution
verifier and a high union ceiling ($0.988$ oracle allocation), the two
undirected budget-scalable baselines catch up at this budget --- the cascade
ties RoR ($0.952$) and random allocation slightly exceeds it ($0.959$) ---
because when almost every query is solvable by \emph{some} draw and verification
is reliable, spreading draws widely eventually succeeds regardless of
\emph{which} model is chosen. RoR's advantage on code is therefore concentrated
at low budget (it matches single-routing at ${\sim}3.2\times$ lower cost) rather
than in a high-budget accuracy ceiling; crucially,
Table~\ref{tab:realverifier-code} shows this low-budget advantage is preserved
under a realistic, imperfect code verifier, whereas the undirected baselines are
not evaluated there because they lack a per-query stopping signal.

\textbf{Distance to the ceiling.} From Table~\ref{tab:results}, oracle
allocation attains $1.0$/$0.944$/$1.0$ at mean cost ${\sim}7$. For the tested
pools, this suggests that knowing \emph{which} cheap model will succeed is
worth more than additional blind budget; the gap between the blue curve and
the amber line can accordingly be read as the value of a better router signal
rather than of more sampling --- the two axes the framework separates.

\begin{figure}[t]\centering
\includegraphics[width=0.9\columnwidth]{figs/fig_verifier.pdf}
\caption{Verifier-quality ablation: policy accuracy at matched mid budget as
the verifier degrades from perfect ($q{=}1$) toward a random pick over drawn
samples. The largest gains (GPQA) degrade fastest as verifier quality drops.}
\label{fig:verifier}
\end{figure}

\textbf{The gains are verifier-gated.} From Fig.~\ref{fig:verifier},
degrading the verifier shrinks the recovered advantage most where it was
largest: at matched mid budget, GPQA falls $0.968\to0.710\to0.675$ as
$q:1.0\to0.8\to0.6$ (MATH-500: $0.887\to0.811\to0.794$; GSM8K:
$0.993\to0.979\to0.970$). At
$q{=}0.6$ the policy still leads all deployable baselines on GPQA ($0.675$ vs
$0.648$ budget-aware best-of-$K$, $0.509$ cascade), but on MATH-500 its margin inverts slightly
against budget-aware best-of-$K$ ($0.794$ vs $0.803$) --- it is speculated
that, when the verifier
cannot be trusted, exploring alternative models buys less than concentrating
draws on one strong model. This mirrors the companion paper's
$\Orepro\le O^{\mathrm{agg}}\le\Oexp$ ordering: test-time
recovery appears available roughly to the extent a deploy-time verifier is.

\begin{figure*}[t]\centering
\includegraphics[width=\textwidth]{figs/fig_pareto_q06.pdf}
\caption{Cost--quality fronts of the budget-scalable policies under a degraded
verifier ($q{=}0.6$; mean over 20 draw orderings). The cascade and random
allocation degrade the most; budget-aware best-of-$K$ becomes the only
competitive baseline, slightly overtaking RoR on MATH-500 while RoR keeps the
lead on GPQA-Diamond.}
\label{fig:pareto_q06}
\end{figure*}

\textbf{Degraded verifier, full front.} Fig.~\ref{fig:pareto_q06} repeats the
whole budget sweep at $q{=}0.6$ rather than a single matched point. Two
observations follow. First, the policies that spread budget across many models
(cascade, random) lose the most --- at matched mid budget they fall to
$0.509$/$0.507$ on GPQA and $0.625$/$0.637$ on MATH-500 --- which is reasonable
because a weak verifier turns every extra candidate into a chance of selecting
a wrong answer. Second, the concentrated policies degrade far less, and the
ordering between RoR and budget-aware best-of-$K$ becomes benchmark-dependent
(RoR ahead on GPQA, $0.675$ vs $0.648$; slightly behind on MATH-500, $0.794$
vs $0.803$), consistent with the interpretation that rerouting pays only when
the verifier can be trusted to recognize a specialist's correct answer.

\textbf{Sensitivity and stability.} Table~\ref{tab:ablations} varies the two
knobs the policy actually has, at a fixed per-query budget $B{=}26$. The prior
pseudo-count $s$ changes accuracy by at most $1.7$~points over an order of
magnitude ($s{=}0.5$ to $8$; GSM8K and MATH-500 are essentially flat), so the
default $s{=}2$ is not tuned to the test set. Calibrating the offline priors
on $30\%$ instead of $70\%$ of the queries costs at most $1.3$~points
(GSM8K $0.983\to0.996$ over the sweep; GPQA $0.735\to0.744$), indicating the
policy needs only a coarse ranking of the models, not precise accuracies.
Across the $20$ draw orderings the standard deviation is
$0.003$/$0.009$/$0.023$ on GSM8K/MATH-500/GPQA respectively --- small relative
to the reported margins, though largest on GPQA, whose test half has only
$99$ queries.

% AUTO-GENERATED by experiments/run_ablations.py -- do not hand-edit.
\begin{table*}[t]\centering\footnotesize
\caption{Sensitivity of the RoR policy at a fixed per-query budget $B=26$ (early stopping spends only part of it, so this is the low-budget operating point of Table~\ref{tab:results}): prior pseudo-count $s$, train fraction used to calibrate offline priors, and stability (mean $\pm$ std over the $20$ draw orderings).}
\label{tab:ablations}
\begin{tabular}{lcccc}\toprule
 & GSM8K & MATH-500 & GPQA & HumanEval+ \\ \midrule
prior $s=0.5$ & 0.989 & 0.828 & 0.744 & 0.897 \\
prior $s=2$ (default) & 0.989 & 0.828 & 0.744 & 0.897 \\
prior $s=8$ & 0.989 & 0.833 & 0.761 & 0.888 \\
\midrule
train frac.\ $0.3$ & 0.983 & 0.827 & 0.735 & 0.914 \\
train frac.\ $0.5$ (default) & 0.989 & 0.828 & 0.744 & 0.897 \\
train frac.\ $0.7$ & 0.996 & 0.843 & 0.744 & 0.906 \\
\midrule
mean $\pm$ std (orderings) & 0.989 $\pm$ 0.003 & 0.828 $\pm$ 0.009 & 0.744 $\pm$ 0.023 & 0.897 $\pm$ 0.011 \\
\bottomrule\end{tabular}\end{table*}


% AUTO-GENERATED by experiments/run_real_costs.py -- do not hand-edit.
\begin{table*}[t]\centering\footnotesize
\caption{Real-price replay (provider-calibrated \$ per M output tokens; exact OpenRouter prices for three listed models, size-tier medians otherwise, accessed July 2026): accuracy at the point nearest mean cost \$0.6 per 1k queries scale. Cost ordering differs from the parameter proxy (e.g.\ Phi-4 is cheaper than most 7--9B models; Llama-3.1-8B is $5\times$ cheaper than same-size peers).}
\label{tab:realcost}
\begin{tabular}{lcccccccc}\toprule
 & \multicolumn{2}{c}{GSM8K} & \multicolumn{2}{c}{MATH-500} & \multicolumn{2}{c}{GPQA} & \multicolumn{2}{c}{HumanEval+} \\
Policy & cost & acc & cost & acc & cost & acc & cost & acc \\ \midrule
\textbf{RoR (ours)} & \textbf{0.06} & \textbf{0.992} & \textbf{0.42} & \textbf{0.903} & \textbf{0.31} & \textbf{0.970} & \textbf{0.25} & \textbf{0.953} \\
budget-aware best-of-$K$ & 0.18 & 0.982 & 0.26 & 0.800 & 0.22 & 0.919 & 0.52 & 0.863 \\
cascade & 0.06 & 0.992 & 0.37 & 0.870 & 0.37 & 0.929 & 0.22 & 0.952 \\
random allocation & 0.23 & 0.992 & 0.61 & 0.880 & 0.59 & 0.887 & 0.41 & 0.960 \\
router (one commit) & 0.46 & 0.976 & 0.46 & 0.784 & 0.14 & 0.566 & 0.46 & 0.817 \\
single-route (best model) & 0.46 & 0.966 & 0.46 & 0.776 & 0.14 & 0.551 & 0.46 & 0.858 \\
\bottomrule\end{tabular}\end{table*}


\textbf{Real-price replay.} The parameter-count proxy preserves size ordering
but not real serving economics, so we replay the whole $q{=}1$ sweep under a
provider-calibrated price vector: the exact OpenRouter output price for the
three pool models listed there, and the median output price of same-size-tier
text models on the same catalog otherwise (7--9B \$0.15/M, $n{=}13$; 14--16B
\$0.24/M; 30--34B \$0.458/M; accessed July 2026).\footnote{\texttt{openrouter.ai/api/v1/models};
prices in \$ per million output tokens.} Real prices decouple cost from size:
Phi-4 (14.7B, \$0.14/M) is cheaper than most 7--9B models, Llama-3.1-8B
(\$0.03/M) is $5\times$ cheaper than same-size peers, and the 32B model costs
only ${\sim}3\times$ the 7--9B median. From Table~\ref{tab:realcost}, the
conclusions carry over under this snapshot: RoR reaches $0.992$ on GSM8K at
\$0.06 (tying the cascade), leads MATH-500 with $0.903$ at \$0.42 against
$0.800$ for best-of-$K$ and $0.870$ for the cascade, and leads GPQA with
$0.970$ at \$0.31 against $0.919$/$0.929$. Notably, budget-aware best-of-$K$
\emph{weakens} under real prices --- committing to one model is harder when
price ordering no longer tracks ability --- while RoR, which reallocates
per query, is less sensitive to the reshuffle. These numbers are specific to
one provider snapshot; the replay itself re-runs in minutes on CPU for any
other price vector.

% AUTO-GENERATED by experiments/run_real_verifier.py -- do not hand-edit.
\begin{table*}[t]\centering\footnotesize
\caption{Agreement-verified replay (consensus threshold $A{=}2$; label-free self-consistency verifier): accuracy (mean spent cost) at two per-query budgets, mean over 20 draw orderings.}
\label{tab:realverifier}
\begin{tabular}{llcccc}\toprule
 & & \multicolumn{2}{c}{$B=26$} & \multicolumn{2}{c}{$B=58$} \\
Benchmark & Policy & acc & cost & acc & cost \\ \midrule
GSM8K & RoR (agreement) & 0.948 & 15.3 & 0.958 & 15.7 \\
 & best-of-$K$ (agreement) & 0.941 & 16.2 & 0.942 & 16.3 \\ \midrule
MATH-500 & RoR (agreement) & 0.757 & 17.3 & 0.789 & 20.7 \\
 & best-of-$K$ (agreement) & 0.779 & 17.4 & 0.787 & 18.9 \\ \midrule
GPQA & RoR (agreement) & 0.552 & 19.8 & 0.536 & 23.1 \\
 & best-of-$K$ (agreement) & 0.572 & 18.5 & 0.569 & 19.0 \\
\bottomrule\end{tabular}\end{table*}


\textbf{A real verifier: agreement-based verification.} Beyond the parametric
$q$, we instantiate a verifier a deployed system can actually run:
\emph{agreement} (self-consistency) --- a drawn sample counts as verified when
its extracted answer matches another drawn sample, the query stops once any
answer reaches a consensus of $A{=}2$ draws, and the final answer is the
plurality cluster. No labels or reward model are involved. From
Table~\ref{tab:realverifier}, the outcome splits by answer space. On GSM8K,
whose free-form numeric answers rarely collide by chance, agreement is
informative: RoR attains $0.958$ against $0.942$ for agreement-gated
best-of-$K$, retaining a small lead. On MATH-500 the two are within a point at
the larger budget ($0.789$ vs $0.787$). On GPQA, a four-option multiple-choice
benchmark, agreement is nearly uninformative --- two wrong draws easily agree
on the same letter --- and the ordering \emph{reverses} ($0.552$ vs $0.572$ at
$B{=}26$), with RoR degrading as budget grows ($0.536$ at $B{=}58$) because
exploring more models manufactures spurious consensus faster. This matches the
parametric prediction that RoR's edge shrinks and can invert in the low-$q$
regime, and yields a concrete deployment rule: agreement-based verification
suffices for open-ended answer spaces, while multiple-choice tasks need a
verifier stronger than consensus.

% AUTO-GENERATED by experiments/run_real_verifier_code.py -- do not hand-edit.
\begin{table*}[t]\centering\footnotesize
\caption{HumanEval+ with a \emph{deployable} code verifier: the in-prompt base tests gate early-stopping (weak, gameable), the full EvalPlus suite is truth. Measured base-verifier false-accept rate 1.0\%; RoR under this real verifier nearly matches its perfect-verifier ceiling (mean over 20 orderings).}
\label{tab:realverifier-code}
\begin{tabular}{lcccc}\toprule
 & \multicolumn{2}{c}{$B=26$} & \multicolumn{2}{c}{$B=58$} \\
Policy & acc & cost & acc & cost \\ \midrule
RoR (base-verifier) & 0.897 & 10.2 & 0.918 & 13.2 \\
RoR (oracle ceiling) & 0.897 & 10.1 & 0.921 & 13.1 \\
best-of-K (base-verifier) & 0.870 & 10.4 & 0.838 & 16.7 \\
\bottomrule\end{tabular}\end{table*}


\textbf{A real verifier for code: partial test suites.} For HumanEval+ the
natural deployable verifier is not agreement but a \emph{partial test suite}:
the original HumanEval base tests are visible in the prompt (weak and gameable),
while the held-out EvalPlus tests are the truth. Using base-test pass as the
early-stopping signal and full-suite pass as ground truth
(Table~\ref{tab:realverifier-code}), the measured false-accept rate is only
$1.0\%$ --- a draw that clears the base tests almost always clears the full
suite. RoR under this real verifier therefore nearly matches its
perfect-verifier ceiling ($0.897$ vs $0.897$ at $B{=}26$; $0.918$ vs $0.921$ at
$B{=}58$) and stays ahead of base-verified best-of-$K$ ($0.870$/$0.838$). This
is the code counterpart to the rule above: where multiple-choice agreement is
too weak to gate recovery, code ships with a reliable partial verifier, so the
test-time recovery RoR exploits is available in practice, not just under an
idealized oracle.

% auto-generated by run_latency.py — mean sequential rounds (round-trips) per query
\begin{table*}[t]\centering\footnotesize
\caption{Latency proxy: mean sequential rounds (round-trips) per query at matched budget. RoR and cascade are adaptive (one draw per round); best-of-$K$ and single-route issue their draws in a single parallel round.}
\label{tab:latency}
\begin{tabular}{lcccc}\toprule
 & GSM8K & MATH-500 & GPQA & HumanEval+ \\ \midrule
single-route & 1.0 & 1.0 & 1.0 & 1.0 \\
best-of-$K$ (parallel) & 1.0 & 1.0 & 1.0 & 1.0 \\
cascade & 1.7 & 3.8 & 3.6 & 2.9 \\
RoR (ours) & 1.2 & 3.6 & 3.3 & 2.5 \\
\bottomrule\end{tabular}\end{table*}


\textbf{Latency.} Cost is not the only deployment axis. RoR and the cascade are
\emph{adaptive} --- each draw's outcome informs the next choice, so they incur a
round-trip per draw --- whereas best-of-$K$ and single-route issue their draws
in a single parallel batch. Table~\ref{tab:latency} reports the mean sequential
rounds per query at matched budget. RoR is sequential ($2.5$ rounds on
HumanEval+, $3.3$ on GPQA), i.e.\ higher latency than the one-round baselines,
but it consistently uses \emph{fewer} rounds than the cascade while reaching at
least its accuracy (HumanEval+ $2.46$ vs $2.87$; GPQA $3.32$ vs $3.59$). The
adaptivity that buys accuracy is paid for in round-trips, and among the adaptive
policies RoR is Pareto-efficient on the accuracy--latency trade-off.

\textbf{UCB variant.} From Table~\ref{tab:results}, the exploration bonus is
at best neutral: it ties the
greedy rule on GSM8K, trails it slightly on MATH-500 ($0.876$ vs $0.887$ at
matched cost), and loses on GPQA at low budgets ($0.603$ vs $0.744$ at mean
cost ${\sim}14$). It is speculated that exploratory draws on a hard benchmark
burn budget that greedy exploitation of the offline prior spends better.
Hence, the greedy rule is recommended as the default.

\section{Discussion}\label{sec:discussion}
\textbf{Verifier dependence.} The policy assumes a usable verifier signal;
where no better-than-random verifier exists, resampling cannot be steered and
the advantage over the budget-aware best-of-$K$ baseline narrows or inverts
(Fig.~\ref{fig:verifier}); the agreement instantiation of
Table~\ref{tab:realverifier} exhibits both regimes on real data. The verifier model is parametric by design --- it
isolates \emph{how much} verifier quality the mechanism needs; plugging in a
concrete process-reward or execution verifier (e.g.\ unit tests for code,
where $q\approx1$ holds naturally) is a direct instantiation rather than a new
method.

\textbf{Cost model.} Costs are proxied by parameter count, which preserves the
pool's cost \emph{ordering} but not exact \$-ratios; re-running the replay with
a deployment's own price/latency vector requires no new generation, only
re-weighting --- Table~\ref{tab:realcost}
demonstrates this with a provider price snapshot. Batched and latency-constrained serving, where draws of one
query can be parallel, is left to future work.

\textbf{Deployment considerations.} The policy's overhead is a few floating-point
operations per step (Section~\ref{sec:method}), so the practical integration
cost lies elsewhere: engineering a verifier signal (a reward model, execution
tests, or agreement across draws), measuring the deployment's own price/latency
vector, and deciding how much per-query latency the sequential draws may add.
When the pool is homogeneous or the benchmark is near-saturated, the simpler
budget-aware best-of-$K$ captures most of the benefit (Table~\ref{tab:results});
the case for full RoR is strongest when models genuinely specialize and a
trustworthy verifier exists.

\textbf{Offline replay.} Replaying pre-computed draws is exact for
correctness-based rewards under the companion protocol's seed-aligned
generation, and makes every number reproducible from the released tensors on
CPU; it does not capture prompt-adaptive behaviors (e.g.\ revising the prompt
between draws).

\section{Conclusion}\label{sec:conclusion}
This work formulated the choice between resampling a committed LLM and
rerouting to an alternative model as a single budgeted decision, proposed an
online allocation policy driven by estimated marginal correctness per unit
cost, and grounded its behavior in the recoverability asymmetry of the
companion analysis. Replay experiments on four regenerated open-model pools
showed that, across these benchmarks and for the tested pools and baselines, the
proposed policy attains a favorable cost--quality Pareto front: it dominates or matches every
budget-scalable baseline and overtakes the single-commit baselines once the
budget affords a second draw, winning chiefly on cost where the benchmark is
saturated and on accuracy where the pool is heterogeneous; the gains are
verifier-gated, as the companion theory predicts. Robustness checks with a
provider-price replay and a label-free agreement verifier show where these
conclusions carry over (real prices; open-ended answer spaces) and where they
weaken (consensus on multiple-choice tasks). In the future, a line of research
is to instantiate stronger learned verifiers (e.g.\ process-reward models) and
to extend the replay protocol to batched and latency-constrained serving.

\bibliographystyle{IEEEtran}
\bibliography{refs}

\end{document}
